\section{Why induction is not circular}\label{ch05:sec:not-circular}
A first reaction to induction is often suspicion: the induction step \emph{assumes} the formula at $k$, so are we assuming what we want to prove? We are not. The induction step proves a conditional statement, ``\emph{if} the formula holds at $k$, \emph{then} it holds at $k+1$,'' and it proves this for every $k$ at once. It never claims that the formula holds at any particular $k$. The base case supplies the first actual truth, and the conditional statement then carries that truth forward one step at a time. Think of an endless row of dominoes. The induction step says that each domino, if it falls, knocks over the next one. The base case pushes the first domino.

Here is a more careful way to see that no natural number can slip through. It is a proof by contradiction, as in Chapter~\ref{ch:proof}. Suppose we have proved the base case $P(0)$ and the induction step, and imagine that $P$ nevertheless fails for some natural number. Among all the natural numbers where $P$ fails, choose the smallest one and call it $m$. Such a smallest number exists because every nonempty set of natural numbers has a smallest member. This is the well-ordering property of Chapter~\ref{ch:proof}, which was stated there for positive integers. Allowing $0$ changes nothing: if $0$ belongs to the set, it is the smallest member, and otherwise the set consists of positive integers. Now look at what we know about $m$:
\begin{itemize}
\item $m\ne0$, because $P(0)$ was proved directly. So $m\ge1$, and $m-1$ is a natural number.
\item $P(m-1)$ is true, because $m-1$ is smaller than $m$, and $m$ is the \emph{smallest} number where $P$ fails.
\item The induction step, applied with $k=m-1$, turns $P(m-1)$ into $P(m)$. So $P(m)$ is true.
\end{itemize}
So $P(m)$ is both false, by the choice of $m$, and true, by the last bullet. This contradiction shows that $P$ fails nowhere: $P(n)$ is true for every natural number $n$. The pattern of this argument, ``choose a smallest counterexample and derive a contradiction,'' is called the \emph{minimal-counterexample} method. Chapter~\ref{ch:proof} used it to prove that every integer is even or odd, and many other proofs begin the same way.

\subsection*{Strong induction}
Sometimes the step to $k+1$ needs more than the case immediately before it. \emph{Strong induction} lets us assume all of $P(0),P(1),\ldots,P(k)$ when proving $P(k+1)$. To prove $P(n)$ for every natural number $n$ by strong induction, prove
\begin{enumerate}
\item the base case: $P(0)$ is true; and
\item the strong induction step: for every natural number $k$, if $P(0),P(1),\ldots,P(k)$ are all true, then $P(k+1)$ is true.
\end{enumerate}
The minimal-counterexample argument shows that this works too. If $m$ were the smallest natural number where $P$ fails, then $m\ne0$ by the base case. Each of $P(0),P(1),\ldots,P(m-1)$ would be true, because each of these numbers is smaller than $m$. The strong induction step with $k=m-1$ would then give $P(m)$, a contradiction. Like ordinary induction, strong induction can start at a number other than $0$.

A classic example shows why the extra strength helps. Recall from Section~\ref{ch02:sec:implication} that an integer $n$ is divisible by an integer $d$ when $n=dk$ for some integer $k$. A positive integer $d$ with this property is called a \emph{divisor} of $n$. An integer $p\ge2$ is \emph{prime} if its only positive divisors are $1$ and $p$. For example, $2$, $3$, $5$ and $7$ are prime, but $12$ is not, since $3$ is a divisor of $12$. By a \emph{product of primes} we mean a product of two or more primes, which need not be different, such as $12=2\cdot2\cdot3$.

\begin{proposition}\label{ch05:prop:prime-product}
Every integer $n\ge2$ is a prime or a product of primes.
\end{proposition}

\begin{proof}
We use strong induction on $n$, starting at $2$. Let $n\ge2$, and assume that every integer $m$ with $2\le m<n$ is a prime or a product of primes. (When $n=2$ there are no integers $m$ with $2\le m<2$, so nothing is assumed. This case plays the role of the base case, and the argument below handles it too, because $2$ is prime.)

If $n$ is prime, there is nothing to prove. Otherwise $n$ has a positive divisor $a$ other than $1$ and $n$, so $n=ab$ for some integer $b$. We check that both $a$ and $b$ lie in the range covered by the assumption, that is, between $2$ and $n-1$:
\begin{itemize}
\item Since $a$ and $n$ are positive, $b=n/a$ is positive; as an integer, $b\ge1$. Hence $n=ab\ge a$. Together with $a\ne1$ and $a\ne n$, this gives $2\le a\le n-1$.
\item Since $a\ge2$, we have $b=n/a\le n/2<n$. Since $a<n$, we have $b=n/a>1$, so $b\ge2$. Hence $2\le b\le n-1$.
\end{itemize}
By the assumption, $a$ is a prime or a product of primes, and so is $b$. Multiplying these together writes $n=ab$ as a product of primes.
\end{proof}

For example, $12=3\cdot4$. The number $3$ is prime and $4=2\cdot2$, so $12=3\cdot2\cdot2$. Ordinary induction would let us assume the result only for $n-1$, and that is useless here: for $n=12$ we need the cases $3$ and $4$, not $11$. We cannot know in advance which smaller numbers $a$ and $b$ will turn up, so we need the result for all of them. Notice also that the proof checked that $a$ and $b$ really are smaller than $n$ and at least $2$ before using the assumption about them. Every use of an induction hypothesis needs a check of this kind. The same situation arises in the proof of Hall's theorem in Chapter~\ref{ch:hall}, which splits a problem into two smaller problems whose sizes are not known in advance.
